% =====================================================================================
% sections/ext_intro.tex: material cut from the Introduction and Related Work on 2026-10-10
% (CUT_PLAN.md, Sec. 4). Extended version only; appendix-ready (\subsection, no \section).
% Where each block came from (all from the pre-cut sections/intro.tex):
%   "Motivation and Costs in Detail"  Introduction, first paragraph (motivating example) and the
%                                     paragraph "Costs" (every cost number of the old intro, with
%                                     its todos; the numbers are those of evaluation.tex,
%                                     generation_quality.tex and v1.tex)
%   "Related Work in Detail"          Related Work, all five paragraphs at full length (published
%                                     costs of zkSNARK provers, per-scheme sampling rules, the
%                                     folklore behind the lower bound, and the long positioning),
%                                     reworded per CUT_PLAN.md decision 4 ("dichotomy" dropped,
%                                     attention "recomputed by the verifier from checked values")
% Labels defined: ext:intro-costs, ext:related (the body does not reference them).
% Labels used (all defined in the body): thm:guess, thm:dich, tab:sok, sec:setup-proof, sec:gen,
%   sec:v1, sec:quality, tab:eval-compare.
% =====================================================================================

\providecommand{\para}[1]{\par\noindent\textbf{#1}\quad}
\providecommand{\todo}[1]{\textbf{[TODO: #1]}}
\providecommand{\sC}{\mathsf{C}}
\providecommand{\Kpre}{\mathsf{K}_{\mathrm{pre}}}

\subsection{Motivation and Costs in Detail}
\label{ext:intro-costs}

\para{Motivation.} A company that pays a cloud provider to serve a model to its users cannot tell
from a response which computation produced it. The provider may serve a cheaper model, or it may
alter the computation on chosen queries to return answers of its choice, which harms any process
that acts on those answers, such as payment approval or fraud flagging. The company cannot re-run
the model to find out, since it lacks the weights, the hardware or both; that is why it outsourced
inference.

\para{Costs.} Beyond the forward pass, the prover's work without V1 grows with the model but not
with the input, and the prover is the protocol's strongest axis. Computing its column openings and
folds with exact int8 tensor-core products instead of float64 makes it 3.7--4.3$\times$ faster on
Llama-2 blocks on an RTX~2080~Ti, and it proves a 256-token GPT-2 response to a 64-token prompt in
0.29\,s. On matched models it is 46--32{,}000$\times$ faster than published zkSNARK provers, partly
because all of them but zkLLM ran on CPUs, and partly because they also prove the operations
without weights, which our verifier recomputes, and many of them hide the weights
(Table~\ref{tab:eval-compare}) \todo{L40S: final range;
the current one is from the pre-improvement L40S run}. The main cost is the proof. It carries the
results of every weight operation and so grows with the prompt: 325\,MB for a 64-token Llama-2-7B
query and 6.6\,GB at 2{,}048 tokens \todo{L40S}, against 183\,kB for zkLLM. V1 makes it
2.0--2.3$\times$ smaller at 2{,}048 tokens with random weights (the full OPT-1.3B: 2.07 to
0.96\,GB) and 1.68--1.77$\times$ on the real smoothed OPT-125M at 512 tokens, at the price of a GKR
in the prover whose work grows with the number of results (OPT-125M at 2{,}048 tokens: 0.50 to
4.85\,s on the RTX~2080~Ti) and a CPU verifier that is 1.6--1.9$\times$ slower at that length
\todo{L40S}. The verifier recomputes attention, so its cost grows with the square of the prompt
length. With eight CPU threads, a client that holds the weights verifies a 64-token Llama-2-7B
query in about 1.5\,s, against 3.5\,s for re-executing the model in fp32, but at 2{,}048 tokens the
two take about the same time, about 85\,s \todo{L40S/EPYC: re-measure with the native attention
kernel}. Our GPU claim decoder and fused exact attention kernel make the streaming GPU verifier
2.7$\times$ faster on OPT-1.3B at that length \todo{L40S, full models}. The protocol is not
zero-knowledge: the claims of $\lceil k/T\rceil$ queries of $T$ tokens reveal a weight matrix with
rows of length $k$. Finally, the provider must serve the integer model. Moving SmoothQuant's
per-channel scales~\cite{xiao2023smoothquant} into the gains of the integer LayerNorm brings its
WikiText-2 perplexity within a factor of 1.03 of fp32 on OPT-125M and OPT-1.3B, from 1.17 and 1.11,
without changing any operation that the protocol checks (Section~\ref{sec:quality})
\todo{OPT-6.7B and Qwen3-4B with smoothing}.

\subsection{Related Work in Detail}
\label{ext:related}

\para{zkSNARK proofs of inference.} A long line of work compiles neural networks into succinct
proofs. For CNNs these include sum-check and GKR provers~\cite{gkr} such as zkCNN~\cite{zkcnn},
zkPyTorch~\cite{zkpytorch} and VerfCNN~\cite{verfcnn}, and pairing- or Plonk-based systems such as
vCNN~\cite{vcnn}, ZEN~\cite{zen}, ZENO~\cite{zeno}, ZKML~\cite{zkml}, EZKL~\cite{ezkl} and
Bionetta~\cite{bionetta}; Mystique~\cite{mystique} is an interactive, VOLE-based alternative. For
language models they include zkLLM~\cite{zkllm}, zkGPT~\cite{zkgpt}, DeepProve~\cite{deepprove},
ZKTorch~\cite{zktorch}, Jolt Atlas~\cite{joltatlas} and the full mode of NanoZK~\cite{nanozk}. They
give computational soundness with proofs of under a kilobyte to tens of megabytes, and many of them
hide the weights. Their provers are far slower than inference. zkLLM takes 620\,s on an A100 for a
2{,}048-token Llama-2-7B prompt, with a 183\,kB proof~\cite{zkllm}; ZKTorch takes 44\,min for one
Llama-2-7B token~\cite{zktorch}; DeepProve takes 34\,s for a 64-token GPT-2 prompt on a 16-core
CPU~\cite{deepprove}. DeepProve and zkAgent~\cite{zkagent} prove a generated sequence in one pass
over the prompt and the output, which is also how we prove generation (Section~\ref{sec:gen}). These
systems prove non-linear operations with lookup arguments, such as zkLLM's tlookup and the
logUp-GKR~\cite{papini2023logupgkr} that DeepProve uses. V1 (Section~\ref{sec:v1}) borrows this
machinery for one narrow purpose, to show that hidden results lie in their requantisation windows.
The weight products still go through Freivalds' check against the existing commitment, and no
activation is committed. Hollow-LLM~\cite{hollowllm} observes that a valid proof over committed
private weights does not establish that the provider performed the computation of a model of the
declared size, since weights with a degenerate algebraic structure keep the declared shapes but are
cheap to evaluate. Our protocol shares this limitation for private weights; for open weights the
commitment is recomputed from the checkpoint. Our protocol gives up the zkSNARKs' small proofs and
weight privacy in exchange for a prover that is orders of magnitude faster.

\para{Spot-checking and sampling.} Anchuri et al.~\cite{anchuri2026verifiable} spot-check random
paths through a committed trace and prove other-model soundness through trace separation.
SLP~\cite{slp} commits every chunk boundary and proves a random subset of chunks with a SNARK; one
audit of its 70B configuration accepts an edited chunk with probability 98.14\%. CommitLLM's routine
audit~\cite{commitllm2026} opens 10 of 80 layers at one token of an audited response.
TensorCommitments~\cite{tensorcommitments} and NanoZK's audit mode~\cite{nanozk} check a subset of
layers, SPEX~\cite{spex} re-executes with a fixed probability, and IMMACULATE~\cite{immaculate}
audits a fraction of requests. TOPLOC~\cite{toploc} and SVIP~\cite{svip} fingerprint hidden states
and target model, prompt or precision substitution. Verde~\cite{verde}, opML~\cite{opml} and
TAO~\cite{tao} are optimistic: an answer stands unless a party that re-runs the model disputes it,
so they rely on an honest party that re-runs every query. Table~\ref{tab:sok} places these designs
under our lower bound. SLP and CommitLLM themselves describe the edit-then-recompute
strategy~\cite{slp,commitllm2026}, and the calculation behind it is classical. It underlies proofs
of retrievability and of data possession~\cite{juels2007por,ateniese2007pdp}, it is why
spot-checkers certify only closeness~\cite{ergun2000spotcheckers}, and it is why PCPs of proximity
guarantee only that their input is close to the language while holographic proofs encode their
input~\cite{bghsv2006robust,bfls1991,chiesa2020fractal}. Theorem~\ref{thm:guess} turns this
calculation into a bound for adaptive, value-aware verifiers that read the whole trace; its tight
case is the optimal search for a hidden object~\cite{chew1967,kadane1971}. Proof-of-Learning~\cite{jia2021pol},
which checks the largest weight updates of a training run, is a precedent for value-dependent
sampling being gamed~\cite{zhang2022pol,fang2023pol}.

\para{Trusted hardware.} Slalom~\cite{slalom} checks the linear layers of a GPU computation from
inside an SGX enclave with Freivalds' algorithm, optionally with a precomputed secret challenge.
TwinShield~\cite{twinshield} and VeriAttn~\cite{veriattn} extend enclave-based checking to
attention, and VeriAttn accepts results within a tolerance. GPU confidential computing on
H100-class GPUs runs the whole model inside an attested enclave on the GPU; a measurement study
finds little overhead inside the GPU and attributes most of the penalty, below 7\% for most LLM
queries, to encrypted transfers over PCIe~\cite{gpucc}. These designs are fast, but their guarantee
rests on the enclave and its vendor. Our protocol needs no trusted hardware, and its exact integer
semantics leaves no tolerance in which to hide a change.

\para{Freivalds-based checks and linear-code commitments.} SafetyNets~\cite{safetynets} uses
interactive proofs for networks with quadratic activations, and ZKML~\cite{zkml} uses Freivalds'
check for linear layers inside a SNARK. Maverick~\cite{maverick} checks every linear layer of a
language model with sparse Freivalds checks against a code-encoded copy of the weights that the
client stores. Gupta et al.~\cite{gupta2026private} run Freivalds' check in the exponent of a
group, against verification data that the model owner generates with secret randomness, for a
client that outsources inference to two non-colluding servers. CommitLLM~\cite{commitllm2026}
checks int8 language models with secret Freivalds vectors precomputed from the weights, at
12--14\% serving overhead, and audits attention without verifying it. LAMP~\cite{lamp} lets a
verifier that holds only a commitment check one matrix product, with Freivalds' check and
Reed--Solomon proximity testing inside a Groth16 SNARK. Our column check is the evaluation phase of
the Ligero and Brakedown commitments~\cite{ligero,brakedown}, and our setup proof
(Section~\ref{sec:setup-proof}) is their testing phase, done once per model in the style of
Ligero. DeepBrake~\cite{deepbrake} consolidates the proximity test and the evaluation check of a
row-wise Reed--Solomon commitment into a single proximity test for evaluation points fixed in
advance, without a sum-check, by fixing the row evaluations before the verifier samples the random
fold. ReedWeave~\cite{reedweave} combines interleaving with FRI-style recursive folding to speed up
Reed--Solomon polynomial commitments. Among zkSNARKs, DeepProve's hash-based configuration commits
with BaseFold~\cite{basefold}. A proximity test of this kind in place of our full-column setup
proof is the natural way to shorten that proof, which we leave to future work.

\para{Where we fit.} We keep the setting of Anchuri et al., a client that holds only a commitment
to the weights, and replace the sampled check, which needs $(1-\varepsilon)N$ paths per layer for
soundness $1-\varepsilon$ (Theorem~\ref{thm:dich}), with a check of every weight operation. Each of
the fully sound protocols above that avoids a SNARK needs more than a cryptographic commitment:
Slalom an enclave, a GPU confidential-computing deployment trusted GPU hardware, Maverick an
encoded copy of the weights, CommitLLM a secret derived from them, and Gupta et al.\ public data
that grows with every matrix and is generated with secret randomness. zkSNARKs that verify against
a Merkle root, such as DeepProve with BaseFold, need no trusted setup but prove far more slowly. To
our knowledge, ours is the first Freivalds-based proof of inference that is fully sound for complete
transformers, with attention recomputed by the verifier from checked values, whose client in
mode~$\sC$ needs no enclave, no stored copy of the weights and no weight-derived secret. It uses no
SNARK, only a hash function and a linear code, and with the setup proof it does not trust the party
that computes the commitment. Its novelty lies in this combination for a weightless,
GPU-less client, not in a new cryptographic primitive: the per-query check is the evaluation
protocol of a linear-code polynomial commitment applied to every matrix product of the model.
